\documentclass[runningheads]{llncs}
\usepackage[T1]{fontenc}
\usepackage{amsmath,amssymb,booktabs,graphicx}
\usepackage[hidelinks]{hyperref}
\hypersetup{pdfauthor={Shayan Bemanian}}
\newcommand{\PublicationResultsRows}{%
LightGCN & $0.1386\pm0.0012$ & $0.0849\pm0.0004$ & $0.1130\pm0.0026$ & $0.0713\pm0.0017$ & $0.1319\pm0.0007$ & $0.0834\pm0.0005$ \\
Full RLMRec & $0.1465\pm0.0023$ & $0.0894\pm0.0010$ & $0.1227\pm0.0015$ & $0.0769\pm0.0010$ & $0.1424\pm0.0009$ & $0.0903\pm0.0005$ \\
Item-only RLMRec & $0.1462\pm0.0013$ & $0.0895\pm0.0006$ & $0.1225\pm0.0016$ & $0.0773\pm0.0005$ & $0.1410\pm0.0006$ & $0.0893\pm0.0002$ \\
SIG & $0.1536\pm0.0016$ & $0.0935\pm0.0007$ & $0.1199\pm0.0010$ & $0.0757\pm0.0006$ & $0.1471\pm0.0002$ & $0.0931\pm0.0002$ \\
}

\newcommand{\CleanGraphMainRows}{%
Amazon & LightGCN & $0.1386\pm0.0012$ & $0.0849\pm0.0004$ \\
Amazon & Full RLMRec & $0.1465\pm0.0023$ & $0.0894\pm0.0010$ \\
Amazon & Item-only RLMRec & $0.1462\pm0.0013$ & $0.0895\pm0.0006$ \\
Amazon & Hybrid framework & $0.1536\pm0.0016$ & $0.0935\pm0.0007$ \\
Yelp & LightGCN & $0.1130\pm0.0026$ & $0.0713\pm0.0017$ \\
Yelp & Full RLMRec & $0.1227\pm0.0015$ & $0.0769\pm0.0010$ \\
Yelp & Item-only RLMRec & $0.1225\pm0.0016$ & $0.0773\pm0.0005$ \\
Yelp & Hybrid framework & $0.1199\pm0.0010$ & $0.0757\pm0.0006$ \\
Steam & LightGCN & $0.1319\pm0.0007$ & $0.0834\pm0.0005$ \\
Steam & Full RLMRec & $0.1424\pm0.0009$ & $0.0903\pm0.0005$ \\
Steam & Item-only RLMRec & $0.1410\pm0.0006$ & $0.0893\pm0.0002$ \\
Steam & Hybrid framework & $0.1471\pm0.0002$ & $0.0931\pm0.0002$ \\
}

\newcommand{\CleanGraphPairedRows}{%
Amazon & $+0.00705$ & $[+0.00295, +0.01115]$ & 0.0354 & 0.0625 \\
Yelp & $-0.00281$ & $[-0.00400, -0.00162]$ & 0.0197 & 0.0625 \\
Steam & $+0.00474$ & $[+0.00376, +0.00572]$ & 0.0017 & 0.0625 \\
}

\newcommand{\PublicationGraphAloneRows}{%
Amazon & 0.13953 & 0.14991 & 0.15518 & $+0.01565$ & $+0.00526$ \\
Yelp & 0.11899 & 0.12383 & 0.10357 & $-0.01542$ & $-0.02026$ \\
Steam & 0.13594 & 0.14201 & 0.14462 & $+0.00869$ & $+0.00262$ \\
}

\newcommand{\PublicationPermutationRows}{%
Amazon & 64D, kd=0.03, 5 seeds, test & 0.16676 & 0.15076 & $-0.01600$ \\
Yelp & 32D, kd=0.1, k=20, seed 2023, validation & 0.12110 & 0.11465 & $-0.00646$ \\
Steam & 32D, kd=0, seed 2023, validation & 0.14462 & 0.13549 & $-0.00914$ \\
}

\title{A Hybrid Framework for Recommendation via LLM Semantic Anchoring and Static Graph Filtering}
\titlerunning{LLM Semantic Anchoring and Static Graph Filtering}
\author{Shayan Bemanian}
\authorrunning{Shayan Bemanian}
\institute{}
\begin{document}
\maketitle

\begin{abstract}
We combine a frozen semantic item graph with item-side contrastive alignment,
using RLMRec's released semantic embeddings.
Over five paired seeds, our framework improves test Recall@20 over full RLMRec
by $0.00705$ on Amazon-book and $0.00474$ on Steam, and loses by $0.00281$
on Yelp. The direction is consistent across all five seeds on each dataset.
Steam inherits Amazon's model hyperparameters, with a different batch size and
stopping budget. Exploratory controls show that removing alignment can retain
a validation advantage over selected baselines, while permuting item labels
on the graph reduces performance. Those controls differ in protocol and do
not isolate the complete mechanism of the main-test advantage. We report the
Yelp loss explicitly; a recorded semantic-agreement prediction failed on
Steam, leaving the cross-dataset boundary unresolved.
\keywords{Recommendation \and Language models \and Graph neural networks \and
Semantic item graphs}
\end{abstract}

\section{Introduction}
RLMRec aligns collaborative representations with semantic vectors obtained
from generated user and item profiles~\cite{rlmrec}. The same item vectors
can also define a similarity graph that supplies information during
propagation, following the frozen-graph idea in FREEDOM~\cite{freedom}.
These uses are complementary: one changes the training objective and the
other changes the representation computed by the model.

Our framework combines a frozen cosine nearest-neighbour item
graph with an item-side alignment loss and a LightGCN collaborative branch.
We compare it with full and item-only RLMRec-Con on Amazon-book, Yelp and
Steam using the same released semantic inputs. Main model comparisons share
a dataset's evaluation and stopping protocol, but the recipes differ in
several architectural and objective choices. Their performance difference
cannot be attributed solely to the addition of an item graph.

The result is a consistent but conditional advantage: the hybrid framework beats
full RLMRec on Amazon and Steam and trails it on Yelp. Steam uses model
hyperparameters transferred from Amazon without a Steam-specific sweep;
its training batch and patience differ from Amazon's. Supporting experiments
remove the alignment loss or permute graph labels to examine how the recipe
uses item semantics. We distinguish these exploratory controls from the
five-seed test comparisons. They motivate further investigation but do not
explain why the recipe loses to RLMRec on Yelp.

\section{Method}
\paragraph{Backbone and objective.}
Let $B$ training triples $(u,i,j)$ contain an observed item $i$ and a sampled
negative $j$. With collaborative embeddings $h$ and $\hat y_{ui}=h_u^\top h_i$,
the averaged BPR loss is
\begin{equation}
 \overline{\mathcal L}_{\rm BPR}=\frac{1}{B}\sum_{(u,i,j)}
 \operatorname{softplus}(\hat y_{uj}-\hat y_{ui}).
\end{equation}
RLMRec-Con adds batch-averaged cosine InfoNCE terms on users and on positive
and negative items, through a shared two-layer projector that maps frozen
1536-dimensional semantic vectors into the collaborative dimension. The full
objective is $\mathcal L_{\rm full}=\overline{\mathcal L}_{\rm BPR}
+\alpha(\overline{\mathcal L}_u+\overline{\mathcal L}_{i^+}
+\overline{\mathcal L}_{i^-})+\beta\|\Theta\|_2^2$, and dropping
$\overline{\mathcal L}_u$ gives the item-only variant.

\paragraph{Frozen semantic item graph.}
The hybrid framework builds a cosine top-$k$ graph $S$ over the item semantic vectors,
excluding self-links, takes the binary symmetric union and normalises it as
$D^{-1/2}SD^{-1/2}$. Item ID embeddings propagate once over this fixed graph and
their unit-normalised output is added to the collaborative item embeddings. The
collaborative branch uses mean layer readout and samples $80\%$ of training
edges per epoch with weights proportional to $1/\sqrt{d_ud_i}$; evaluation uses
the complete training graph. Alignment, when enabled, applies only to positive
items' fused representations. Setting $\alpha=0$ removes the alignment term
entirely and leaves the graph as the only semantic input.

\paragraph{Relation to prior recipes.}
Semantic and learned item graphs are established: LATTICE learns item relations
from side features~\cite{lattice}, FREEDOM freezes them and samples the
collaborative graph~\cite{freedom}, and SAGERec constructs an LLM-derived item
graph for sequential recommendation~\cite{sagerec}. AlphaRec shows that item
language representations alone support strong recommendation~\cite{alpharec}.
Our contribution is not the idea of a semantic item graph but the comparison of complete recipes against profile alignment on shared
inputs and a matched main-comparison protocol within each dataset.

We audited our implementation against the released FREEDOM code and record the
differences rather than claiming a reproduction. Ours excludes self-links where
upstream leaves them eligible, symmetrises before normalising where upstream applies a
directed row normalisation, unit-normalises the structural output before
addition, prunes $20\%$ of edges where upstream drops $80$--$90\%$, and replaces
upstream's feature-level BPR terms with an InfoNCE term on a single frozen
semantic modality. The resulting framework is FREEDOM-inspired, not
FREEDOM with different features, and our results do not speak to the performance
of the original model.

\section{Experimental Design}
\paragraph{Data and protocol.}
We use RLMRec's released interaction matrices and semantic embeddings for all
three benchmarks: Amazon-book (11,000 users, 9,332 items, 120,464 training
interactions), Yelp (11,091 / 11,010 / 166,620) and Steam (23,310 / 5,237 /
316,190). Splits are approximately 60/20/20 and we verified that the three
matrices share no nonzero pair on any dataset, and that every semantic array has
the expected entity count and 1,536 finite dimensions. Collaborative adjacency
uses training data only. All-ranking evaluation masks training positives;
validation positives remain eligible test candidates. Evaluation
metrics are macro-averaged over users with positives in the evaluated split.

We cannot verify a complete split-filtered generation history for the released
profiles. The upstream description cites GPT-3.5-turbo and text-embedding-ada-002,
but published examples are not generation logs. Our results are conditional on
these fixed inputs; disjoint interaction splits do not certify the absence of
contamination upstream.

\paragraph{Configurations.}
All compared models use 32-dimensional collaborative embeddings, three
collaborative layers, Adam at learning rate $10^{-3}$ and no optimiser weight
decay. RLMRec variants use $\alpha=0.01$, temperature $0.2$, keep rate $0.8$ and
$\beta=10^{-7}$. Batch size is 1024 on Amazon and 4096 on Yelp and Steam, fixed
within each dataset; cross-dataset behaviour is therefore not an isolated effect
of dataset identity. The hybrid framework uses $\alpha=0.03,k=10$ on Amazon and
$\alpha=0.1,k=20$ on Yelp, with one item-graph layer and pruning ratio $0.2$.

Steam inherits Amazon's model hyperparameters, including $\alpha$, $k$,
embedding dimension, layer counts, temperature, regularization and pruning.
Its batch size and stopping budget differ. The recorded study transferred
this recipe without a Steam-specific model hyperparameter sweep; each
confirmation run evaluates the test set once at its best validation checkpoint.
Earlier Steam development runs and historical tests remain in the artifact.
This is retrospective transfer evidence, not a claim of an untouched test set
or a fully preregistered study.

\paragraph{Stopping and checkpoint selection.}
Validation Recall@20 is evaluated every three epochs. Patience is five
non-improving checks for the Amazon and Yelp main comparisons and 25 for
Steam, matched across the four models within each dataset. The best checkpoint
is restored before one final test evaluation per run. All 60 manifest runs
ended by early stopping. The paired seeds are 1, 7, 13, 42 and 2024.

Separate Yelp diagnostics at RLMRec weight $\alpha=0.03$ show that increasing
patience from five to 20 permits recovery from an early low-score phase.
They are validation-only follow-ups and do not replace the main comparisons
at $\alpha=0.01$. This finding cautions that a matched stopping rule need not
be equally suitable for different optimization trajectories. We have not
repeated all main comparisons under alternative patience values.

\paragraph{Uncertainty and repeatability.}
We report sample standard deviations and pointwise paired 95\% Student-$t$
intervals over five seeds. Holm adjustment covers twelve contrasts per
metric: full and item-only versus LightGCN, and item-only and the hybrid framework
versus full, on each of the three datasets. We also report exhaustive two-sided
sign-flip values; their minimum with five pairs is $0.0625$.

Four repeated executions per RLMRec variant at Amazon seed 42 give Recall@20
SDs of $0.000536$ for full and $0.000344$ for item-only. This establishes
execution variation at that particular setting, including variation when the
selected epoch is unchanged. It does not establish a universal detection
threshold or quantify noise on Yelp and Steam. Main uncertainty is assessed
with the paired seed intervals, not a threshold extrapolated from these repeats.

\section{Results}
\begin{table}[t]
\caption{Test Recall and NDCG at 20, mean $\pm$ sample standard deviation over
five paired seeds. Batch size is matched within each dataset. Steam model hyperparameters transfer from Amazon; batch size and patience differ.}
\label{tab:main}
\centering\small
\begin{tabular}{llrr}
\toprule Dataset & Model & Recall@20 & NDCG@20\\\midrule
\CleanGraphMainRows
\bottomrule\end{tabular}
\end{table}

\paragraph{The graph recipe against profile alignment.}
Table~\ref{tab:effects} gives the paired differences. The hybrid framework exceeds full
RLMRec by $0.00705$ Recall@20 on Amazon (4.8\% relative) and $0.00474$ on Steam
(3.3\% relative), and falls below it by $0.00281$ on Yelp (2.3\% relative).
The direction holds in all five seeds per dataset. All three paired $t$
comparisons survive the stated Holm adjustment; the exact sign-flip values
are $0.0625$, so significance depends on the parametric analysis.

The Steam margin is evidence that the Amazon model recipe can transfer to
a further benchmark without a new model hyperparameter sweep. It does not
remove the unequal development histories: Yelp received validation tuning,
Amazon was the development benchmark, and Steam uses a longer stopping budget.

\begin{table}[t]
\caption{Hybrid framework minus full RLMRec, test Recall@20 over five seeds.
Pointwise paired $t$ intervals; $p_h$ is Holm-adjusted within the twelve-contrast
family, and $p_s$ is the exact two-sided sign-flip value.}
\label{tab:effects}
\centering\small
\begin{tabular}{lrrrr}
\toprule Dataset & Difference & 95\% $t$ interval & $p_h$ & $p_s$\\\midrule
\CleanGraphPairedRows
\bottomrule\end{tabular}
\end{table}

\paragraph{Removing alignment: exploratory evidence.}
Table~\ref{tab:graphalone} reports development-seed validation scores after
setting $\alpha=0$. On Amazon and Steam the recipe exceeds the listed
LightGCN and item-only baselines, whereas on Yelp it falls below both.
These are single-seed diagnostics, not five-seed test results or comparisons
with full RLMRec. They also have unequal patience on Amazon and Yelp:
Amazon uses 5/5/25 checks for LightGCN/item-only/graph, Yelp 25/5/5, and
Steam 25/25/25. Thus only Steam matches patience across these rows.
Other recipe differences remain even when alignment is removed, so the
contrasts do not isolate a causal effect of graph structure alone.

\begin{table}[t]
\caption{Validation Recall@20 on development seed 2023 with alignment switched
off ($\alpha=0$), against LightGCN and item-only RLMRec. Batch size is
1024 on Amazon and 4096 on Yelp and Steam; patience differs as stated in text.}
\label{tab:graphalone}
\centering\small
\begin{tabular}{lrrrrr}
\toprule
Dataset & LightGCN & Item-only & Graph, $\alpha=0$ & $-$LightGCN & $-$Item-only\\
\midrule
\PublicationGraphAloneRows
\bottomrule
\end{tabular}
\end{table}

\paragraph{Graph-label permutation controls.}
An isomorphic relabelling keeps the graph topology and degree multiset while
changing the correspondence between items and neighbourhoods.
Table~\ref{tab:permutation} shows lower scores after this intervention on all
three datasets. The Amazon row is a five-seed 64D test comparison with legacy
provenance limitations; Yelp and Steam are single-seed 32D validation
comparisons. These rows do not form one matched three-dataset experiment.

On Steam, without alignment, the permuted recipe reaches $0.13549$ compared
with $0.14462$ for the real graph and $0.13594$ for LightGCN. The real and
permuted runs select epoch indices 357 and 324. This supports the relevance
of item-to-neighbourhood correspondence in that setting. It does not prove
that all of the main-test gain is semantic, rule out optimization effects,
or explain the Yelp comparison with RLMRec.

\begin{table}[t]
\caption{Isomorphic graph permutation: identical topology and degree multiset,
reassigned item labels. Difference is permuted minus real. Settings differ
across rows; the Amazon baseline has incomplete historical protocol metadata.}
\label{tab:permutation}
\centering\small
\resizebox{\textwidth}{!}{\begin{tabular}{llrrr}
\toprule
Dataset & Setting & Real graph & Permuted & Difference\\
\midrule
\PublicationPermutationRows
\bottomrule
\end{tabular}}
\end{table}

\section{Limitations and Conclusion}
The Yelp loss is unexplained. We measured the agreement between each item's
semantic neighbourhood and its collaborative co-occurrence, expecting it to
predict where the graph helps, and recorded a prediction for Steam before its development runs. The prediction failed. A subsequent correction to that diagnostic's
null distribution produced an ordering consistent with the outcomes, but
selecting a measure after seeing three results is not evidence, and we do not
present it as a mechanism. The artifact contains the original prediction, its
falsification, and the correction.

Our scope is one collaborative backbone, three released benchmarks, fixed
offline semantic inputs and Apple MPS execution, for which fixed-seed execution variation was measured at one Amazon setting. Batch size and patience vary across datasets, so cross-dataset contrasts
are not isolated effects of dataset identity. The Amazon and Yelp configurations have unequal tuning histories; the Steam model hyperparameters transfer without a new sweep, but
the training protocol changes. We do not
measure profile-generation cost, and removing user alignment does not remove
user interaction data or address cold start, since item profiles may themselves
be derived from user-authored text.

The supported conclusion is that the proposed hybrid framework improves
five-seed test performance over full RLMRec on Amazon-book and Steam and
reduces it on Yelp under the recorded protocols. Alignment-removal and
permutation diagnostics support further study of semantic graph correspondence,
but their scope and mismatched settings prevent a complete causal account.
The failed semantic-agreement prediction leaves the Yelp loss unresolved.

\paragraph{Artifact and provenance.} The accompanying artifact includes all
60 main results, their hash-checked manifest, per-seed statistics, graph
controls, the recorded Steam prediction, the implementation audit and the
scripts that regenerate the tables. It preserves the exact source snapshot
for the separate alignment-weight study; complete historical source bytes
are not recovered for every other run. Raw inputs are acquired separately
from RLMRec, and their upstream generation history is incompletely verified.

\paragraph{AI assistance.} AI assistants supported framing, code, statistical
checks, experiment execution, literature search and drafting. All reported
numbers come from archived training outputs and the artifact's analysis scripts.

\clearpage
\begin{thebibliography}{8}
\bibitem{rlmrec} Ren, X., Wei, W., Xia, L., Su, L., Cheng, S., Wang, J., Yin, D.,
Huang, C.: Representation Learning with Large Language Models for Recommendation.
In: The Web Conference (2024). \doi{10.1145/3589334.3645458}
\bibitem{freedom} Zhou, X., Shen, Z.: A Tale of Two Graphs: Freezing and Denoising
Graph Structures for Multimodal Recommendation. In: ACM Multimedia (2023).
\doi{10.1145/3581783.3611943}
\bibitem{lightgcn} He, X., Deng, K., Wang, X., Li, Y., Zhang, Y., Wang, M.:
LightGCN: Simplifying and Powering Graph Convolution Network for Recommendation.
In: SIGIR (2020). \doi{10.1145/3397271.3401063}
\bibitem{alpharec} Sheng, L., Zhang, A., Zhang, Y., Chen, Y., Wang, X., Chua, T.-S.:
Language Representations Can Be What Recommenders Need: Findings and Potentials.
In: ICLR (2025). \url{https://openreview.net/forum?id=eIJfOIMN9z}
\bibitem{lattice} Zhang, J., Zhu, Y., Liu, Q., Wu, S., Wang, S., Wang, L.:
Mining Latent Structures for Multimedia Recommendation. In: ACM Multimedia (2021).
\doi{10.1145/3474085.3475259}
\bibitem{sagerec} Cui, W., Lam, H.-K.: SAGERec: Semantic-Aware Global Graph-Enhanced Representation
Learning for Sequential Recommendation. Electronics \textbf{14}(24), 4844 (2025).
\doi{10.3390/electronics14244844}
\end{thebibliography}
\end{document}
